% 数论（综述）
% license CCBY4
% type Wiki

本文根据 CC-BY-SA 协议转载翻译自维基百科\href{https://en.wikipedia.org/wiki/Number_theory}{相关文章}

\begin{figure}[ht]
\centering
\includegraphics[width=6cm]{./figures/ae528936f6668ed9.png}
\caption{} \label{fig_SHuLun_1}
\end{figure}
数论是纯数学的一个分支，主要致力于研究整数以及算术函数。数论学家研究素数，以及由整数构造的数学对象（例如有理数）的性质，或作为整数的推广而定义的对象（例如代数整数）。

整数既可以单独研究，也可以作为某些方程的解来研究（丢番图几何）。数论中的问题常常能够通过研究某些解析对象来理解，例如黎曼ζ函数，它以某种方式编码了整数、素数或其他数论对象的性质（解析数论）。另外，还可以研究实数与有理数之间的关系，例如研究无理数如何用分数来逼近（丢番图逼近）。

数论与几何一样，是数学最古老的分支之一。数论的一个特点是，它处理的往往是非常容易理解却极难解决的命题。例如，费马大定理在最初提出后的 358 年才被证明，而哥德巴赫猜想自 18 世纪提出以来至今仍未解决。德国数学家卡尔·弗里德里希·高斯曾经说过：“数学是科学的女王，而数论是数学的女王。”[1] 数论一度被视为纯粹数学的典范，认为它在数学之外没有任何应用，直到 20 世纪 70 年代，人们才发现素数可被用作构建公钥密码算法的基础。
